%% ============================================================
%% 分片 10：原文第 4 节「Compactification And The Geometry Of Hitchin's Moduli Space」（原文 glshort.tex 起于第 2325 行）
%% 由 tools/harvmac2tex.py 从 arXiv 源码（hep-th/0604151v3）自动生成。
%% 只译英文散文；公式、编号、\label、\cite 一律不改。
%% ============================================================

\subsection{Complex Structures Of $\MH$}

\label{sec:complexstr}

Our next task is to understand more concretely the complex
structures on $\MH$.

\par\medskip{\it Description By Holomorphic Data}

Let us look at the hyper-Kahler quotient of $\EUW$ from the
standpoint of complex structure $I$. The combination
$\nu_I=\mu_J+i\mu_K$ is holomorphic in this complex structure:
% 式 (4.20)
\setcounter{equation}{19}
\begin{equation}\label{eq:uttup}
\nu_I=-{1\over \pi}\int_C|d^2z|\,\Tr\,\epsilon\,D_{\bar
z}\phi_z.
\end{equation}
 $\nu_I$ is holomorphic in complex structure $I$ because
it is the moment map derived from $\Omega_I$, which is of type
$(2,0)$ in complex structure $I$ (explicitly $\nu_I$ is
holomorphic because it depends only on $A_{\bar z}$ and $\phi_z$).
In such problems, it is often convenient \cite{ref61}\ to consider
separately the vanishing of the holomorphic moment map $\nu_I$ and
the real moment map $\mu_I$.

Any  connection $A$ on a smooth $G$-bundle $E$ over a Riemann
surface $C$ automatically turns $E$ into a holomorphic $G$-bundle
(which we will denote by the same name, hoping this causes no
confusion). One simply defines the $\bar\partial$ operator as
$\bar D=d\bar z\,D_{\bar z}$, which in complex dimension one
trivially obeys $\bar D^2=0$. The vanishing of the holomorphic
moment map $\nu_I$ tells us that $\bar D\phi=0$; in other words,
$\varphi=\phi_z\,dz$ is a holomorphic section of $K_C\otimes {\rm
ad}(E)$, where $K_C$ is the canonical line bundle of $C$.
Differently put, $\varphi$ is a holomorphic one-form valued in the
adjoint representation, that is in the adjoint bundle of $E$. We
will call a pair $(E,\varphi)$ consisting of a $G$-bundle $E$ and
a holomorphic section $\varphi$ of $K_C\otimes {\rm ad}(E)$ a
Hitchin pair.  $\varphi$ is often called the Higgs field, and the
bundle $E$ endowed with the choice of $\varphi$ is also called a
Higgs bundle.

To obtain the moduli space $\MH$, we must also set to zero the
real moment map $\mu_I$ and divide by the group of $G$-valued
gauge transformations.  However, as exploited in \cite{ref20}\ (and as
is often the case in moduli problems) there is a simpler way to
understand these combined steps.  The space $\EUW_0$ of zeroes of
the holomorphic moment map $\nu_I$ admits the action not just of
ordinary $G$-valued gauge transformations, but of gauge
transformations valued in the complexification $G_{\Bbb{C}}$ of
$G$.  This is manifest from the holomorphy of $\nu_I$. Thus  we
can perform on $\EUW_0$ either of the following two operations:

{\it (I)} Restrict to the subspace of $\EUW_0$ with $\mu_I=0$ and
divide by $G$-valued gauge transformations.

 {\it (II)} Divide $\EUW_0$ by $G_{\Bbb{C}}$-valued gauge
transformations.

 Operation {\it (I)} gives the desired moduli space
$\MH$, but operation {\it (II)} is much easier to understand and
nearly coincides with it.  (The reason they nearly coincide is
that almost every orbit of complex gauge transformations contains
a unique orbit of ordinary gauge transformations on which
$\mu_I=0$.)  The easiest way to understand operation {\it (I)} is
often to first understand operation {\it (II)} and then understand
its relation to operation {\it (I)}.

In the present case, the result of operation {\it (II)} is easy to
describe. It means that we do not care about the particular choice
of $A_{\bar z}, $ $\phi_z$, but only about the holomorphic
$G$-bundle $E$ that is determined by $A_{\bar z}$, and the
associated holomorphic section $\varphi$ of $K_C\otimes {\ad}(E)$.
Thus, in the present example, operation {\it (II)} gives us the
set of equivalence classes of Hitchin pairs $(E,\varphi)$.

This set is for many purposes a very good approximation to $\MH$;
for example, they differ in rather  high codimension if $C$ has
genus at least 2. To be more precise (most of the present paper
does not depend on the details),  we need the notion of stability.
For $G=SU(2)$, we interpret $E$ as a rank two holomorphic vector
bundle over $C$, and $\varphi$ as a holomorphic map $\varphi:E\to
E\otimes K_C$. A line bundle ${\cal L}\subset E$ is called
$\varphi$-invariant if $\varphi({\cal L})\subset {\cal L}\otimes
K_C$. A Hitchin pair $(E,\varphi)$ is called stable if every
$\varphi$-invariant line bundle ${\cal L}\subset E$ has negative
first Chern class. It is called semistable if each such ${\cal L}$
has non-positive first Chern class. For general $G$, one must
consider $\varphi$-invariant reductions of the structure group of
$E$ to a maximal parabolic subgroup $P$ of $G_{\Bbb{C}}$. The
bundle $E$ with such a reduction has a natural first Chern class
(because $P$ has a distinguished $U(1)$), and the pair
$(E,\varphi)$ is called stable (or semistable) if for every such
reduction the first Chern class is negative (or nonpositive). A
pair that is semistable but not stable is called strictly
semistable.

Stability is a mild condition in the sense that, for example, if a
pair $(E,\varphi)$ is stable, then so is every nearby pair.
Moreover, the pairs that are not stable have high codimension (if
the genus of $C$ is greater than 1).

Now we can go back to the question of comparing operations {\it
(I)} and {\it (II)}, or equivalently, describing in holomorphic
terms the moduli space $\MH$ of solutions of Hitchin's equations.
The result proved in \cite{ref20}\ is that $\MH$ is the same as the
``moduli space of stable pairs $(E,\varphi)$,'' i.e., stable
Hitchin pairs. We get this moduli space by throwing away unstable
Hitchin pairs, imposing a certain equivalence relation on the
semistable ones, and then dividing by the complex-valued gauge
transformations.  This slightly modified   version of operation
{\it (II)} -- call it operation {\it (II)$'$} -- agrees precisely
with operation {\it (I)}.

In sum, $\MH$ parametrizes the Hitchin pairs $(E,\varphi)$ that
are stable, as well as certain equivalence classes of strictly
semistable Hitchin pairs.  The stable pairs correspond to smooth
points in $\MH$ as well as (for some $G$) certain orbifold
singularities. The strictly semistable pairs generally lead to
singularities in $\MH$ that are worse than orbifold singularities.

\par\medskip{\it Analog For Complex Structure $J$}

All this has an analog in complex structure $J$. The holomorphic
moment map in complex structure $J$ is $\nu_J=\mu_K+i\mu_I$:
% 式 (4.21)
\setcounter{equation}{20}
\begin{equation}\label{eq:nuJ}
\nu_J =-{i\over 2\pi}\int_C\,\Tr\,\epsilon{\cal F}
\end{equation}
 To
construct $\MH$, we first impose the vanishing of $\nu_J$. A zero
of $\nu_J$ is simply a pair $A,\phi$ such that the curvature
${\cal F}$ of the $G_{\Bbb{C}}$-valued connection $\CA=A+i\phi$ is
equal to zero. If therefore we were to divide the zero set of
$\nu_J$ by the group of $G_{\Bbb{C}}$-valued gauge
transformations, we would simply get the set ${\cal Y}_0$ of all
homomorphisms $\vartheta:\pi_1(C)\to G_{\Bbb{C}}$, up to
conjugation. This is operation {\it (II)}.

Instead, to get $\MH$, we are supposed to carry out operation {\it
(I)}.  In other words, we are supposed to set $\mu_J=0$, i.e., to
impose the condition $D^*\phi=0$, and divide only by the
$G$-valued gauge transformations.  As  in complex structure $I$,
in comparing these two operations, a notion of stability intrudes.
For $G=SU(2)$, a homomorphism $\vartheta:\pi_1(C)\to G_{\Bbb{C}}$
is considered stable if the monodromies cannot be simultaneously
reduced to the triangular form % 式 (4.22)
\setcounter{equation}{21}
\begin{equation}\label{eq:yto}
\left(\begin{matrix} \alpha &
\beta \\ 0 & \alpha^{-1}\end{matrix}\right).
\end{equation}
   If $\vartheta$ has such a
reduction, we call it strictly semistable.  (In complex structure
$J$, there is no notion of an unstable representation.) We
consider two strictly semistable representations to be equivalent
if they have the same diagonal monodromy elements, that is, the
same $\alpha$'s.  Note that each such equivalence class has a
distinguished representative with $\beta=0$.

For any $G$, the analog of putting the monodromies in triangular
form is to conjugate them into a parabolic subgroup $P$ of $G$.  A
representation is stable if it cannot be so conjugated, and
otherwise strictly semistable.  The equivalence relation on
strictly semistable representations has a natural analog for any
$G$ (two strictly semistable representations that both reduce to
$P$ are equivalent if the two flat $P$ bundles become equivalent
when projected to the maximal reductive subgroup of $P$).

A theorem of Corlette \cite{ref58}\ and Donaldson  (see the appendix
to \cite{ref20}) is that $\MH$, in complex structure $J$, is  the
moduli space ${{\cal Y}}$ that parametrizes stable homomorphisms
$\vartheta:\pi_1(C)\to G_{\Bbb{C}}$, as well as the equivalence
classes of semistable ones.

\par\medskip{\it Analog For $I_w$}

What we have just said has a direct analog in each of the complex
structure $I_w$, $w\not= 0,\infty$.  The holomorphic  variables in
complex structure $I_w$ are $A_{\bar z}-w\phi_{\bar z}$,
$A_z+w^{-1}\phi_z$.  Two of Hitchin's equations combine to the
holomorphic condition $[D_{\bar z}-w\phi_{\bar
z},D_z+w^{-1}\phi_z]=0$, and the third is a moment map condition.
The holomorphic condition says that the complex-valued connection
with components $(A_{\bar z}-w\phi_{\bar z},A_z+w^{-1}\phi_z)$ is
flat.  Setting to zero the moment map and dividing by gauge
transformations gives the same moduli space ${{\cal Y}}$ that we
found in complex structure $J$.  Thus all the complex structures
$I_w$, $w\not=0,\infty$ are equivalent.

In fact, in complex structure $I=I_{w=0}$, we can define a
symmetry group ${\cal U}\cong \Bbb{C}^*$ that acts on a Higgs
bundle by $(E,\varphi)\to (E,\lambda\varphi)$,
$\lambda\in\Bbb{C}^*$. (The stability condition is invariant under
this transformation.) For $|\lambda|=1$, this reduces to the group
${\cal U}_1$, described in section 第~\ref{sec:mhhyper}~节, which is a symmetry
group of the hyper-Kahler metric of $\MH$. For $|\lambda|\not=1$,
we do not get a symmetry of the hyper-Kahler metric of $\MH$, but
as is shown on pp. 107-8 of \cite{ref20}, we get a group of
diffeomorphisms of $\MH$ that preserves complex structure $I$ and
transforms the other complex structures by % 式 (4.23)
\setcounter{equation}{22}
\begin{equation}\label{eq:trymodo}
I_w\to
I_{\lambda^{-1}w},
\end{equation}
 generalizing eqn. \eqref{eq:zungo}\ to
$|\lambda|\not=1$.

\subsection{Hitchin's Fibrations}

\label{sec:hitchfib}

Now we will discuss in more detail the geometry of $\MH$.

If $E$ is a stable $G$-bundle, then the pair $(E,0)$ is a stable
Hitchin pair.  So we get a natural embedding of $\M$, the moduli
space of stable $G$-bundles on $C$, into the Hitchin moduli space
$\MH$.  $\M$ is a holomorphic submanifold of $\MH$ in complex
structure $I$, since it is defined by the equations $\phi_z=0$,
which are holomorphic in complex structure $I$.  In complex
structures $J$ and $K$, $\M$ is not holomorphic. Instead, it is
Lagrangian, since $\phi=0$ implies $\delta\phi=0$ and hence
$\mu_J=\mu_K=0$.

If $E$ is stable, then the Hitchin pair $(E,\varphi)$ is stable
for every $\varphi\in H^0(C,K_C\otimes\ad(E))$.  Moreover, the
tangent bundle to ${\EUM}$ at the point $E$ is $H^1(C,\ad(E))$
(which parametrizes first order deformations of $E$), so by Serre
duality, the cotangent space to $\M$ at the point $E$ is
$H^0(C,K_C\otimes\ad(E))$. So $\varphi$ takes values in this
cotangent space, and the space of all $(E,\varphi)$ with stable
$E$ is the cotangent bundle $T^*\M$. We thus actually get an
embedding of $T^*\M$ in $\MH$.  The holomorphic symplectic form
$\Omega_I$ of $\MH$ in complex structure $I$ restricts on $T^*\M$
to its natural symplectic structure as a holomorphic cotangent
bundle (that is, for any local holomorphic coordinates $q^\alpha$
on $\M$, we have $\Omega_I=\sum_\alpha dp_\alpha\wedge dq^\alpha$
for some holomorphic functions $p_\alpha$ that are linear on the
fibers of $T^*\M\to \M$).\footnote{This statement is true because the
analogous statement is actually true on the space $\EUW$ of all
pairs $(A,\phi)$ before taking the hyper-Kahler quotient.  $\EUW$
can be understood as the cotangent space of the space ${\EUA}$ of
connections, with $A$ parameterizing ${\EUA}$ and $\phi$
parameterizing the cotangent space to ${\EUA}$; $\Omega_I$ is the
natural holomorphic symplectic form of ${\EUW}$ regarded as
$T^*{\EUA}$.  In general, this type of structure persists in
taking a complex symplectic quotient (in the present case, setting
$\nu_I=0$ and dividing by $G_{\Bbb{C}}$-valued gauge
transformations).}

The image of $T^*\M$ in $\MH$ is not all of $\MH$ because a
Hitchin pair $(E,\varphi)$ may be stable even if $E$ is unstable.
However, the stable Hitchin pairs $(E,\varphi)$ for which $E$ is
unstable are a set of very high codimension. Upon throwing away
this set, $\MH$ becomes isomorphic to $T^*\M$, and has a natural
map to $\M$ by forgetting $\varphi$. Even though it is only
defined away from a set of very high codimension, this map is
extremely useful.  We will call it Hitchin's second fibration (the
first one being another map that we introduce presently).

\par\medskip{\it The Hitchin Fibration}

Another natural operation in complex structure $I$, apart from
mapping $\varphi$ to zero, is to calculate the gauge-invariant
polynomials in $\varphi$. For $G=SU(2)$, this simply means that we
consider the quadratic Casimir operator $w=\Tr\,\varphi^2$. Since
$\bar D\varphi=0$, we have $\bar\partial w=0$, so $w$ is a
holomorphic quadratic differential, taking values in
${\EUBB}=H^0(C,K_C^2)\cong \Bbb{C}^{3g-3}$. The Hitchin fibration,
as it is most commonly called, is the map $\pi:\MH\to {\EUBB}$
obtained by mapping the pair $(E,\varphi)$ to $w=\Tr\,\varphi^2$.

For any $G$, the Hitchin fibration is defined similarly, except
that one characterizes $\varphi$ by all of its independent
Casimirs (that is, all of the independent $G$-invariant
polynomials in $\varphi$), not just the quadratic one. For
example, for $G=SU(N)$, we  define $w_n=\Tr\,\varphi^n$,
$n=2,\dots,N$, and let ${\EUBB}=\oplus_{n=2}^NH^0(C,K_C^n)$. The
Hitchin fibration is then the map that takes $(E,\varphi)$ to the
point $(w_2,w_3,\dots,w_n)\in {\EUBB}$. For any $G$, one considers
instead of $\Tr\,\varphi^n$ the appropriate independent
homogeneous $G$-invariant polynomials ${\cal P}_i$. The number of
these polynomials equals the rank $r$ of $G$, and their degrees
$d_i$ obey the identity % 式 (4.24)
\setcounter{equation}{23}
\begin{equation}\label{eq:yty}
\sum_{i=1}^r(2d_i-1)={\rm
dim}(G).
\end{equation}
 For example, for $G=SU(N)$, the $d_i$ are $2,3,\dots,N$,
whence $\sum_i(2d_i-1)=N^2-1={\rm dim}(G)$.  And finally in
general we define ${\EUBB}=\oplus_i H^0(C,K_C^{d_i})$.

\def\CMF{{\cmmib F}}
Since $\dim \,H^0(C,K_C^n)=(2n-1)(g-1)$, it follows from \eqref{eq:yty}\
that the complex dimension of ${\EUBB}$ is $(g-1){\rm dim}(G)$,
which equals the dimension of $\M$, and  one half of the dimension
of $\M_H$. The Hitchin fibration $\pi:\MH\to {\EUBB}$ is
surjective, as we will discuss momentarily.   As the base
${\EUBB}$ of the Hitchin fibration $\pi:\MH\to {\EUBB}$ has one
half the dimension of $\MH$, it follows that the dimension of a
typical fiber $\CMF$ of $\pi$ is also half the dimension of $\MH$
and so equal to the dimension of ${\EUBB}$:
% 式 (4.25)
\setcounter{equation}{24}
\begin{equation}\label{eq:yrof}
\dim\,\CMF=\dim\,{\EUBB}={1\over
2}\dim\,\MH=(g-1)\dim\,G.
\end{equation}

Let us explain qualitatively why the Hitchin fibration is
surjective. For example, take $G=SU(2)$. Pick a stable $SU(2)$
bundle $E$. Consider the equations $\Tr\,\varphi^2=w$, where
$\varphi $ varies in the $(3g-3)$-dimensional space
$H^0(C,K_C\otimes \ad(E))$ and $w$ is a fixed element of the
$(3g-3)$-dimensional space ${\EUBB}=H^0(C,K_C^2)$. This is a
system of $3g-3$ quadratic equations for $3g-3$ complex variables.
The number of solutions is generically $2^{3g-3}$.  A similar
counting can be made for other $G$.

\par\medskip{\it Complete Integrability}

 
 

 We now want to explain one of Hitchin's main
results \cite{ref62}: $\MH$ is a completely integrable Hamiltonian system
in the complex structure $I$.  We can find ${1\over 2}{\rm
dim}\MH$ functions $H_a$ on $\MH$ that are holomorphic in complex
structure $I$, are algebraically independent, and are
Poisson-commuting using the Poisson brackets obtained from the
holomorphic symplectic form $\Omega_I$.\footnote{The reader may be
unaccustomed to completely integrable systems in this holomorphic
sense.  {}From such systems, one can extract completely integrable
Hamiltonian systems in the ordinary real sense (and moreover,
interesting and significant constructions arise in this way; see
\cite{ref63,ref64} for some examples). We pick $C$ to admit
a real structure -- that is, an involution that reverses its
complex structure. This induces real structures on $\M$ and $\MH$.
By specializing to a real slice in $\MH$, one gets then completely
integrable Hamiltonian systems in which the phase space
coordinates, the symplectic structure, and the commuting
Hamiltonians are all real.}

In fact, we can take the $H_a$ to be linear functions on
${\EUBB}$, since the dimension of ${\EUBB}$ is the same as the
desired number of functions.   We will explain the construction
first for $G=SU(2)$. We begin by picking  a basis $\alpha_a$,
$a=1,\dots,3g-3$ of the $(3g-3)$-dimensional space $ H^1(C,T_C)$,
which is dual to $H^0(C, K_C^2)\cong {\EUBB}$. (Here $T_C$ is the
holomorphic tangent bundle to $C$.) We represent $\alpha_a$ by
$(0,1)$-forms valued in $T_C$, which we call by the same name, and
we define % 式 (4.26)
\setcounter{equation}{25}
\begin{equation}\label{eq:miro}
H_a=\int_C\alpha_a \wedge \Tr\,\varphi^2.
\end{equation}
 We
claim that these functions are Poisson-commuting with respect to
the holomorphic symplectic form $\Omega_I$.

A natural proof uses the fact that the definition of the $H_a$
makes sense on the infinite-dimensional space $\EUW$, before
taking the hyper-Kahler quotient.  Using the symplectic structure
$\Omega_I$ on $\EUW$ to define Poisson brackets, the $H_a$ are
obviously Poisson-commuting. For in these Poisson brackets, given
the form \eqref{eq:ninto}\ of $\Omega_I$, $\varphi_z$ is conjugate to
$A_{\bar z}$ and commutes with itself. But the $H_a$ are functions
of $\varphi_z$ only, not $A_{\bar z}$.

The functions $H_a$ can be restricted to the locus with $\nu_I=0$,
and then, because they are invariant under the
$G_{\Bbb{C}}$-valued gauge transformations, they descend to
holomorphic functions on $\MH$. A general property of symplectic
reduction (in which one sets to zero a moment map, in this case
$\nu_I$, and then divides by the corresponding group, in this case
the group of $G_{\Bbb{C}}$-valued gauge transformations) is that
it maps Poisson-commuting functions to Poisson-commuting
functions.  So the $H_a$ are Poisson-commuting as functions on
$\MH$.  There are enough of them to establish the complete
integrability of $\MH$.

One point to note in this construction is that the $H_a$, as
functions on $\EUW$, depend on the particular choice of $T$-valued
$(0,1)$-forms $\alpha_a$ that represent the cohomology classes.
Any choice will do, but we have to make a choice.  But after
restricting and descending to $\MH$, the functions we get on $\MH$
depend only on the cohomology classes of the $\alpha_a$. In fact,
once we have $\bar D\varphi=0$ and hence
$\bar\partial\Tr\,\varphi^2=0$, the $H_a$ are clearly invariant
under $\alpha_a\to\alpha_a+\bar\partial\epsilon_a$.

The Poisson-commuting functions $H_a$ generate commuting flows on
$\MH$ that are holomorphic in complex structure $I$.  Complex tori
admit commuting flows, and one might surmise that the orbits
generated by the $H_a$ are complex tori at least generically. This
follows from general results about holomorphic symplectic
structures and compactness of the fibers of the Hitchin fibration
and can also be demonstrated  more directly, using the technique
of the spectral cover \cite{ref20,ref62}. This technique has
many applications in further development of this subject, as will
be explained elsewhere \cite{ref26}. In this paper, we will get as far
as we can without using the spectral cover construction.

The analog of the above construction for any $G$ is to replace
$\Tr\,\varphi^2$ by a general gauge-invariant  polynomial ${\cal
P}_i$ of degree $d_i$. The associated commuting Hamiltonians take
the form $H_{\alpha,i}=\int_C \alpha \,{\cal P}_i(\varphi)$, for
$\alpha\in H^1(C,K_C^{1-d_i})$. A simple dimension counting, using
the dimension formula \eqref{eq:yty}, shows that these Hamiltonians are
precisely sufficient in number to establish the complete
integrability of $\MH$.  By using the fact that the fibers of the
Hitchin fibration are compact, so that a holomorphic function must
be a pullback from the base of this fibration, one can show that
the commuting Hamiltonians generate the ring of holomorphic
functions on $\MH$.

One easy and important consequence of complete integrability is
that the fibers of the Hitchin fibration are Lagrangian
submanifolds in the holomorphic symplectic structure $\Omega_I$ or
equivalently in the real symplectic structures $\omega_J$ and
$\omega_K$.  Indeed, a fiber of this fibration is defined by
equations $H_k-h_k=0$, where $H_k$ are the commuting Hamiltonians
and $h_k$ are complex constants.  In general, the zero set of a
middle-dimensional collection of Poisson-commuting functions, such
as $H_k-h_k$ in the present case, is Lagrangian.
